\documentclass[12pt]{article}

%% notes in the margin
\usepackage{marginnote}
%% verbatim JavaScript
\usepackage[pdftex]{insdljs}
%% hyperref as usual
\usepackage{hyperref}

%% patching a definition in hpdftex.def, which ignores the backgroundcolor attribute on PushButtons even though pdf supports it
\makeatletter
  \def\PDFForm@Push{%
    /Subtype/Widget%
    \Fld@annotflags
    \Fld@annotnames
    /FT/Btn%
    \Fld@flags
    /H/P%
    /BS<</W \Fld@borderwidth/S/\Fld@borderstyle>>%
    \ifcase0\ifnum\Fld@rotation=\z@   \else 1\fi
            \ifx\Fld@bordercolor\relax\else 1\fi
            \ifx\Fld@bcolor\relax\else 1\fi %% FR added this line
            \space
    \else
      /MK<<%
        \ifnum\Fld@rotation=\z@
        \else
          /R \Fld@rotation
        \fi
        \ifx\Fld@bordercolor\relax
        \else
          /BC[\Fld@bordercolor]%
        \fi
        \ifx\Fld@bcolor\relax %% FR added this line
        \else                 %% FR added this line
          /BG[\Fld@bcolor]%   %% FR added this line
        \fi                   %% FR added this line
      >>%
    \fi
    /A<</S/JavaScript/JS(\Hy@escapestring{\Fld@onclick@code})>>%
    \Fld@additionalactions
  }%

%% library JavaScript functions
\begin{insDLJS}{mysetup}{My JavaScript Setup}
function toggle(name) {
  var f = this.getField(name);
  if (color.equal(f.fillColor,color.transparent))
    f.fillColor=color.white;
  else
    f.fillColor=color.transparent;
}
\end{insDLJS}

\begin{document}

% all form fields must be in a single global Form environment
\begin{Form}

% no bordercolor amounts to making the border transparent so that it is invisible
% print is the default default
\def\DefaultOptionsofPushButton{print,bordercolor=}

%\PushButton[name={test2},onclick={app.alert(this.getField('test').display)}]{dsfasfs}
%\PushButton[onclick={this.getField('test').lineWidth=0; this.getField('test').strokeColor=color.transparent}]{dsfasfs}
%$a\sqrt{\PushButton[name={test},onclick={this.getField('test').userName='testtest'}]{\ensuremath{\frac{\bullet}{\sqrt{2}}}}}\sin x$

%% create a button '#1Control' with showing #2 and a button '#1' in the margin showing #3
%% clicking on the former toggles the latter
\newcommand{\button}[3]{\PushButton[name=#1Control,backgroundcolor=,onclick={toggle('#1')}]{\ensuremath{#2}}\marginnote{\PushButton[name=#1,backgroundcolor=1 1 1]{\ensuremath{#3}}}}

Clicking on the bound variable $x$ toggles the visibility of its type $\mathit{nat}$, which is displayed in the margin.

%% a formula with a button
\[\lambda \button{test}{x}{\mathit{nat}}.x+1\]

\end{Form}

\end{document}